\documentclass[a4paper, serif, 10pt]{moderncv}

%% moderncv
% style and color
\moderncvstyle{banking}
\moderncvcolor{black}

% renew moderncv command to adapt for CJK colon
\renewcommand*{\cvitem}[3][.25em]{%
  \ifstrempty{#2}{}{\hintstyle{#2}：}{#3}%
  \par\addvspace{#1}}

\renewcommand*{\cvitemwithcomment}[4][.25em]{%
  \savebox{\cvitemwithcommentmainbox}{\ifstrempty{#2}{}{\hintstyle{#2}：}#3}%
  \setlength{\cvitemwithcommentmainlength}{\widthof{\usebox{\cvitemwithcommentmainbox}}}%
  \setlength{\cvitemwithcommentcommentlength}{\maincolumnwidth-\separatorcolumnwidth-\cvitemwithcommentmainlength}%
  \begin{minipage}[t]{\cvitemwithcommentmainlength}\usebox{\cvitemwithcommentmainbox}\end{minipage}%
  \hfill% fill of \separatorcolumnwidth
  \begin{minipage}[t]{\cvitemwithcommentcommentlength}\raggedleft\small\itshape#4\end{minipage}%
  \par\addvspace{#1}}

%% page layout/margins
\usepackage[top=2.5cm, bottom=2.5cm, left=1.5cm, right=1.5cm]{geometry}

\nopagenumbers{}

%% Babel config for English language
\usepackage[english]{babel}

%% fontspec
\usepackage{fontspec}

\IfFontExistsTF{Linux Libertine}{
  \setmainfont[Ligatures={TeX, Common}, Numbers=Lining, ItalicFont=Linux Libertine]{Linux Libertine}
}{}
\IfFontExistsTF{Linux Libertine O}{
  \setmainfont[Ligatures={TeX, Common}, Numbers=Lining, ItalicFont=Linux Libertine O]{Linux Libertine O}
}{}

%% CTeX
% CJK support, used to show CJK characters in the resume
%
% - fontset=none: disable builtin fontset but instead set the CJK font manually
% - heading=false: disable ctex heading
% - punct=kaiming: use kaiming punctuations styles for CJK
% - scheme=plain: use plain scheme, do not override \normalsize font size
% - space=auto: space settings for CJK characters
%
% ref:
% - http://ctan.mirrorcatalogs.com/language/chinese/ctex/ctex.pdf
\usepackage[UTF8, heading=false, punct=kaiming, scheme=plain, space=auto]{ctex}

\IfFontExistsTF{Noto Serif CJK SC}{
  \setCJKmainfont{Noto Serif CJK SC}
}{}
\IfFontExistsTF{Noto Sans CJK SC}{
  \setCJKsansfont{Noto Sans CJK SC}
}{}

%% line spacing
\usepackage{setspace}
\setstretch{1.125}

%% URL styling - use normal text instead of monospace
\urlstyle{same}

% Auto-underline all links
% Defer until \begin{document} so hyperref is already loaded
\AtBeginDocument{%
  \let\oldhref\href
  \renewcommand{\href}[2]{\oldhref{#1}{\underline{#2}}}%
}

%% Patch \httplink and \httpslink to handle full URLs
% The original moderncv macros always prepend http:// or https:// to URLs.
% This causes issues when the URL already has a protocol (e.g., https://example.com)
% resulting in malformed URLs like https://https://example.com.
% This patch checks if the URL already contains "://" and uses it directly if so.
% Note: We use \str_set:Nx (x-expansion) to expand macros like \@homepage first.
\ExplSyntaxOn
\str_new:N \g_yamlresume_url_str

\RenewDocumentCommand{\httpslink}{O{}m}{
  \str_set:Nx \g_yamlresume_url_str {#2}
  \str_if_in:NnTF \g_yamlresume_url_str {://}
    { \url{#2} }
    { \url{https://#2} }
}

\RenewDocumentCommand{\httplink}{O{}m}{
  \str_set:Nx \g_yamlresume_url_str {#2}
  \str_if_in:NnTF \g_yamlresume_url_str {://}
    { \url{#2} }
    { \url{http://#2} }
}
\ExplSyntaxOff

\name{鈴木 彩花}{}
\title{顧客価値と事業成長を両立するプロダクトマネージャー}
\phone[mobile]{+81 90-1234-5678}
\email{ayaka.suzuki@example.jp}
\homepage{https://ayakasuzuki.example.jp}

\address{日本、東京都、千代田区、東京都千代田区丸の内1-2-3、100-0005}{}{}

\extrainfo{{\small \faLinkedin}\ \href{https://www.linkedin.com/in/ayaka-suzuki}{@ayaka-suzuki} {} {} {} • {} {} {}
{\small \faGithub}\ \href{https://github.com/ayaka-suzuki}{@ayaka-suzuki}}

\begin{document}

\maketitle

\section{基本情報}

\cvline{}{ユーザー中心のプロダクトマネジメントを専門とし、SaaSプロダクトの企画から改善まで一貫してリードしてきたプロダクトマネージャーです。データに基づく仮説立案と定量・定性調査を組み合わせ、KPI改善と顧客満足度向上を両立させます。ステークホルダーとの合意形成やクロスファンクショナルなチーム推進を強みとし、グローバル展開を見据えたプロダクト戦略にも経験があります。}
\section{学歴}

\cventry{2014年4月～2018年3月}
        {学士、経営工学、成績：3.7}
        {早稲田大学}
        {\url{https://www.waseda.jp}}
        {}
        {\begin{itemize}
\item 経営工学を専攻し、データ分析・マーケティング・組織マネジメントを学ぶ
\item チームでのプロジェクト演習を通じて、課題発見から解決策の提案までを経験
\item 卒業研究では「顧客ロイヤルティ向上要因の分析」をテーマに調査・分析を実施
\end{itemize}
\textbf{コース}：マーケティングリサーチ、データサイエンス入門、プロジェクトマネジメント、組織行動論、サービスイノベーション}
\section{職歴}

\cventry{2021年4月～現在}
        {シニアプロダクトマネージャー}
        {株式会社ネクストソリューションズ}
        {\url{https://www.next-sol.example.jp}}
        {}
        {\begin{itemize}
\item B2B SaaSプロダクト「TeamFlow」のプロダクトマネージャーとして、ロードマップ策定から機能開発、ローンチまでを統括
\item ユーザーインタビューや利用ログ分析に基づき、オンボーディング機能を再設計し、アクティベーション率を45\%から68\%に改善
\item 営業・カスタマーサクセス・開発チームと連携し、顧客フィードバックを優先順位付けする仕組みを構築
\item 四半期ごとのKPIレビューを主導し、プロダクト改善サイクルを定着させた
\end{itemize}
\textbf{キーワード}：プロダクト戦略、ロードマップ、ステークホルダー管理、KPI改善、データ分析}

\cventry{2018年4月～2021年3月}
        {プロダクトマネージャー}
        {株式会社デジタルイノベーション}
        {\url{https://www.digital-innovation.example.jp}}
        {}
        {\begin{itemize}
\item 自社メディア向けCMSのリニューアルプロジェクトに参画し、要件定義からリリースまでを主導
\item 編集者・営業チームへのヒアリングを実施し、業務効率化に貢献する編集ワークフロー機能を企画・実装
\item A/Bテストを活用したUI改善を繰り返し、記事編集時間を平均30\%削減
\item エンジニア・デザイナー約10名のチームをプロジェクトマネジメントとして支援
\end{itemize}
\textbf{キーワード}：要件定義、アジャイル開発、A/Bテスト、UI改善、プロジェクト推進}
\section{言語}

\cvline{日本語}{ネイティブ・母国語レベル \hfill \textbf{キーワード}：ビジネス文書作成、プレゼンテーション}
\cvline{英語}{完全な実務レベル \hfill \textbf{キーワード}：TOEIC 905、海外チームとの会議}
\section{スキル}

\cvline{プロダクトマネジメント}{エキスパート \hfill \textbf{キーワード}：ロードマップ策定、要求分析、KPI設計、顧客インタビュー、プライオリティ管理}
\cvline{データ分析}{上級 \hfill \textbf{キーワード}：SQL、Google Analytics、Tableau、A/Bテスト、コホート分析}
\cvline{UI/UXデザイン}{中級 \hfill \textbf{キーワード}：Figma、プロトタイピング、ユーザビリティテスト、ワイヤーフレーム}
\section{受賞・表彰}

\cventry{2024年3月}
        {株式会社ネクストソリューションズ}
        {2023年度 ベストプロダクト賞}
        {}
        {}
        {\begin{itemize}
\item TeamFlowの大型リニューアルプロジェクトをリードし、顧客満足度の大幅向上と売上成長に貢献したことが評価され受賞
\end{itemize}}
\section{資格・免状}

\cventry{2022年6月}
        {PMI日本支部}
        {PMP}
        {\url{https://www.pmi.org/certifications/project-management-pmp}}
        {}
        {}
\section{出版・著作}

\cventry{2023年5月}
        {データドリブンなプロダクト改善の実践}
        {note}
        {\url{https://note.example.jp/ayaka-suzuki/datadriven-pm}}
        {}
        {\begin{itemize}
\item プロダクトマネージャー向けに、KPI設計やA/Bテストの進め方、ステークホルダーを巻き込む方法を事例付きで紹介
\end{itemize}}
\section{プロジェクト}

\cventry{2022年7月～2023年2月}
        {B2B SaaSプロダクトにおけるオンボーディング機能のリニューアルプロジェクト}
        {TeamFlow オンボーディング改善}
        {\url{https://www.next-sol.example.jp/teamflow}}
        {}
        {\begin{itemize}
\item ユーザーインタビューと利用ログ分析を実施し、離脱要因を特定
\item ステップ形式のガイドとテンプレート機能を新たに設計・実装
\item アクティベーション率を45\%から68\%に改善し、解約率を削減
\end{itemize}
\textbf{キーワード}：オンボーディング、ユーザーリサーチ、プロダクト改善}

\cventry{2019年9月～2020年3月}
        {自社CMSに搭載する編集ワークフロー機能の企画・設計プロジェクト}
        {メディア向けCMS編集ワークフロー設計}
        {\url{https://www.digital-innovation.example.jp/cms}}
        {}
        {\begin{itemize}
\item 編集者と営業チームの業務プロセスを可視化し、課題を整理
\item 承認フローとタスク管理機能の要件定義を実施
\item 編集時間を平均30\%削減し、業務効率化に貢献
\end{itemize}
\textbf{キーワード}：業務分析、要件定義、UX設計}
\section{興味・関心}

\cvline{旅行}{国内旅行、写真}
\cvline{読書}{ビジネス書、デザイン}
\section{ボランティア活動}

\cventry{2022年6月～現在}
        {メンター}
        {PMラーニングコミュニティ}
        {\url{https://www.pm-learning.example.jp}}
        {}
        {\begin{itemize}
\item 若手プロダクトマネージャー向けの勉強会で講師を務め、ロードマップ策定やKPI設計の実践的なノウハウを共有
\item 参加者からの相談に定期的に応じ、キャリア形成の支援を行った
\end{itemize}}

\end{document}